\documentclass[aspectratio=169,11pt]{beamer}
\usepackage{../../../shared/theme/esmad_beamer_theme}
\usepackage{amsmath,amssymb}
\usepackage{booktabs}
\usepackage{array}

\title[Causal Econometrics]{Causal Econometrics}
\subtitle{From Prediction to Identification}
\date{\today}

\newcommand{\E}{\mathbb{E}}
\newcommand{\Var}{\mathrm{Var}}
\newcommand{\Cov}{\mathrm{Cov}}
\newcommand{\ATE}{\mathrm{ATE}}
\newcommand{\ATT}{\mathrm{ATT}}
\newcommand{\plim}{\operatorname{plim}}

\begin{document}

\begin{frame}
\titlepage
\end{frame}

\begin{frame}{Learning Goals}
\begin{itemize}
\item distinguish prediction, association, and causal identification;
\item define causal estimands before selecting an estimator;
\item explain how selection, confounding, simultaneity, and omitted variables create endogeneity;
\item connect panel data, Difference-in-Differences, instrumental variables, and counterfactual methods to their identifying assumptions;
\item evaluate diagnostics, uncertainty, falsification tests, and limits of causal claims.
\end{itemize}
\end{frame}

\begin{frame}{Outline}
\tableofcontents
\end{frame}

\section{Prediction, Association, and Causality}

\begin{frame}{Three Different Questions}
\begin{columns}[T]
\begin{column}{0.32\textwidth}
\begin{block}{Prediction}
What value of \(Y\) should we expect for a new observation with features \(X\)?
\[
\widehat{Y}=f(X)
\]
Accuracy is judged out of sample.
\end{block}
\end{column}
\begin{column}{0.32\textwidth}
\begin{block}{Association}
How does the conditional distribution of \(Y\) vary with \(X\)?
\[
\E[Y\mid X=x]
\]
The relation may reflect many mechanisms.
\end{block}
\end{column}
\begin{column}{0.32\textwidth}
\begin{alertblock}{Causality}
How would \(Y\) change if treatment \(D\) were set differently?
\[
\E[Y(1)-Y(0)]
\]
This requires a counterfactual argument.
\end{alertblock}
\end{column}
\end{columns}
\end{frame}

\begin{frame}{Econometrics Starts with the Data-Generating Process}
Consider
\[
Y_i=\beta D_i+\gamma X_i+u_i.
\]
If
\[
\E[u_i\mid D_i,X_i]=0,
\]
the coefficient on \(D_i\) can have a causal interpretation under the stated structural assumptions.

If instead
\[
\Cov(D_i,u_i)\neq 0,
\]
ordinary least squares mixes the effect of \(D_i\) with mechanisms contained in \(u_i\).

\begin{alertblock}{Identification question}
What variation in \(D\) is sufficiently exogenous to reveal the causal effect we want?
\end{alertblock}
\end{frame}

\section{Estimands and Counterfactuals}

\begin{frame}{Define the Estimand First}
For unit \(i\), let \(Y_i(1)\) and \(Y_i(0)\) denote potential outcomes:
\[
\tau_i=Y_i(1)-Y_i(0).
\]
Two common population estimands are
\[
\ATE=\E[Y(1)-Y(0)]
\]
and
\[
\ATT=\E[Y(1)-Y(0)\mid D=1].
\]

\begin{block}{Practical consequence}
ATE and ATT answer different questions. A method that identifies one does not automatically identify the other.
\end{block}
\end{frame}

\begin{frame}{The Fundamental Missing-Data Problem}
The observed outcome is
\[
Y_i=D_iY_i(1)+(1-D_i)Y_i(0).
\]
For each unit, one potential outcome is observed and the other is missing.

\begin{columns}
\begin{column}{0.48\textwidth}
\begin{block}{Observed}
Treated units reveal \(Y(1)\).\\
Untreated units reveal \(Y(0)\).
\end{block}
\end{column}
\begin{column}{0.48\textwidth}
\begin{alertblock}{Unobserved}
Treated units hide \(Y(0)\).\\
Untreated units hide \(Y(1)\).
\end{alertblock}
\end{column}
\end{columns}

\vspace{0.3cm}
Causal inference is a disciplined strategy for constructing credible counterfactual comparisons.
\end{frame}

\section{Endogeneity and Bias}

\begin{frame}{Where Endogeneity Comes From}
\begin{itemize}
\item \textbf{Omitted variables}: an unobserved factor affects treatment and outcome.
\item \textbf{Simultaneity}: treatment and outcome are jointly determined.
\item \textbf{Measurement error}: the regressor is observed with error.
\item \textbf{Selection}: treatment is chosen partly because of expected outcomes.
\item \textbf{Dynamic feedback}: past outcomes influence later treatment.
\end{itemize}

\begin{alertblock}{Econometric symptom}
The problematic component of treatment is correlated with the regression disturbance.
\end{alertblock}
\end{frame}

\begin{frame}{Omitted-Variable Bias}
Suppose the true model is
\[
Y=\beta D+\gamma Z+u,
\]
but \(Z\) is omitted. The simple-regression coefficient converges to
\[
\plim\widehat{\beta}
=
\beta+\gamma\frac{\Cov(D,Z)}{\Var(D)}.
\]

Bias therefore requires both:
\begin{enumerate}
\item \(Z\) affects the outcome;
\item \(Z\) is associated with treatment.
\end{enumerate}

\begin{block}{Important}
Adding controls is useful only when the controls correspond to a defensible causal adjustment strategy.
\end{block}
\end{frame}

\section{Identification Strategies}

\begin{frame}{A Map of Identification Strategies}
\small
\begin{tabular}{p{2.9cm}p{4.1cm}p{5.2cm}}
\toprule
\textbf{Design} & \textbf{Identifying variation} & \textbf{Core assumption} \\
\midrule
Randomized experiment & Random assignment & Randomization and design conditions \\
Selection on observables & Covariate-adjusted comparison & Exchangeability + overlap \\
Panel fixed effects & Within-unit change & Time-invariant heterogeneity removable \\
Difference-in-Differences & Differential post-treatment change & Parallel untreated trends \\
Instrumental variables & Exogenous treatment shift & Relevance + exclusion + independence \\
Regression discontinuity & Assignment threshold & Continuity near cutoff \\
Synthetic control & Weighted donor trajectory & Credible pre-treatment fit \\
\bottomrule
\end{tabular}
\end{frame}

\begin{frame}{Panel Data and Fixed Effects}
A canonical model is
\[
Y_{it}=\beta D_{it}+\alpha_i+\lambda_t+u_{it},
\]
where \(\alpha_i\) captures time-invariant unit heterogeneity and \(\lambda_t\) common time shocks.

The within transformation removes \(\alpha_i\):
\[
Y_{it}-\bar{Y}_i
=
\beta(D_{it}-\bar{D}_i)
+
(u_{it}-\bar{u}_i).
\]

\begin{alertblock}{What fixed effects do not solve}
Time-varying confounders, reverse causality, anticipation, and endogenous treatment timing remain possible.
\end{alertblock}
\end{frame}

\begin{frame}{Difference-in-Differences}
For treated group \(T\) and control group \(C\),
\[
\widehat{\tau}_{DID}
=
(\bar{Y}_{T,post}-\bar{Y}_{T,pre})
-
(\bar{Y}_{C,post}-\bar{Y}_{C,pre}).
\]

\begin{block}{Counterfactual logic}
The control group's change estimates how the treated group would have evolved without treatment.
\end{block}

\begin{alertblock}{Central assumption}
Without treatment, treated and control groups would have followed parallel trends over the relevant horizon.
\end{alertblock}
\end{frame}

\begin{frame}{Event Studies}
A common specification is
\[
Y_{it}
=
\alpha_i+\lambda_t+
\sum_{k\neq -1}\beta_k\mathbb{1}\{t-T_i=k\}+u_{it}.
\]

\begin{itemize}
\item pre-treatment coefficients help reveal whether groups were already diverging;
\item post-treatment coefficients describe dynamic effects relative to the reference period;
\item heterogeneous effects and staggered adoption require care with two-way fixed-effects estimators.
\end{itemize}

\begin{block}{Do not over-interpret pre-trend tests}
Failure to reject pre-trends is not proof of parallel trends; power may be weak.
\end{block}
\end{frame}

\begin{frame}{Instrumental Variables and 2SLS}
Suppose \(D\) is endogenous and \(Z\) is a candidate instrument.

First stage:
\[
D_i=\pi_0+\pi_1Z_i+\pi_2X_i+v_i.
\]

Second stage:
\[
Y_i=\beta_0+\beta_1\widehat{D}_i+\beta_2X_i+\varepsilon_i.
\]

\begin{columns}
\begin{column}{0.48\textwidth}
\begin{block}{Relevance}
\(Z\) must move treatment.
\end{block}
\end{column}
\begin{column}{0.48\textwidth}
\begin{alertblock}{Exclusion}
\(Z\) must affect \(Y\) only through the relevant treatment channel.
\end{alertblock}
\end{column}
\end{columns}
\end{frame}

\begin{frame}{What Does IV Identify?}
With heterogeneous treatment effects and standard IV assumptions, the estimand may be
\[
LATE=\E[Y(1)-Y(0)\mid\text{compliers}].
\]

\begin{itemize}
\item the effect can be local to units whose treatment changes because of the instrument;
\item weak instruments produce unstable estimates;
\item a strong first stage does not validate the exclusion restriction.
\end{itemize}

\begin{alertblock}{Interpretation matters}
IV is not a generic repair for endogeneity. The scientific meaning of the instrument determines the meaning of the estimate.
\end{alertblock}
\end{frame}

\begin{frame}{Synthetic Control}
For one treated unit and donor units \(j=1,\ldots,J\),
\[
\widehat{Y}_{1t}(0)=\sum_{j=1}^{J}w_jY_{jt},
\qquad
w_j\geq 0,\quad \sum_jw_j=1.
\]
The post-treatment effect is
\[
\widehat{\tau}_t=Y_{1t}-\widehat{Y}_{1t}(0).
\]

\begin{itemize}
\item pre-treatment fit is central;
\item donor support determines credibility;
\item placebo units and placebo dates contextualize the observed gap.
\end{itemize}
\end{frame}

\section{Observed Confounding and Robust Estimation}

\begin{frame}{Selection on Observables}
Let
\[
e(X)=P(D=1\mid X)
\]
be the propensity score.

Under
\[
(Y(1),Y(0))\perp D\mid X,
\]
plus positivity and consistency, treatment effects can be identified by appropriate adjustment.

\begin{alertblock}{Key limitation}
Propensity-score machinery cannot adjust for an important unmeasured confounder without additional identifying structure.
\end{alertblock}
\end{frame}

\begin{frame}{Weighting and Overlap}
For the ATE, inverse-probability weights use
\[
w_i=\frac{D_i}{e(X_i)}+\frac{1-D_i}{1-e(X_i)}.
\]

When \(e(X)\) approaches 0 or 1, weights become extreme.

\begin{itemize}
\item inspect propensity-score distributions;
\item inspect effective sample size and weight concentration;
\item define trimming or truncation transparently;
\item report when the effective target population changes.
\end{itemize}
\end{frame}

\begin{frame}{Doubly Robust Logic}
Augmented weighting combines outcome regressions
\[
\widehat{\mu}_1(X),\qquad \widehat{\mu}_0(X)
\]
with a propensity model \(\widehat e(X)\).

\begin{block}{Double robustness}
Under standard conditions, consistency can survive misspecification of one nuisance model, not failure of the causal identification assumptions.
\end{block}

The phrase \emph{doubly robust} concerns estimator nuisance models; it does not make an observational design assumption-free.
\end{frame}

\section{Diagnostics and Uncertainty}

\begin{frame}{Diagnostics Depend on the Design}
\begin{itemize}
\item \textbf{Panel models}: residual dependence, clustering level, within-unit variation.
\item \textbf{DiD}: pre-treatment dynamics, placebo dates, comparison-group composition.
\item \textbf{IV}: first-stage strength, reduced form, instrument interpretation.
\item \textbf{Synthetic control}: pre-treatment fit, donor weights, placebo gaps.
\item \textbf{Weighting/matching}: overlap, balance, weight concentration.
\end{itemize}

\begin{alertblock}{General principle}
Diagnostics can falsify some designs. They rarely prove that an untestable identifying assumption is true.
\end{alertblock}
\end{frame}

\begin{frame}{Sampling Uncertainty Is Not Identification Uncertainty}
A narrow confidence interval is conditional on the maintained design and model assumptions.

It does not quantify uncertainty about:
\begin{itemize}
\item an invalid exclusion restriction;
\item an unmeasured confounder;
\item anticipation before treatment;
\item interference between units;
\item a structurally poor donor pool.
\end{itemize}

\begin{block}{Reporting standard}
Separate sampling uncertainty from sensitivity to assumptions and model specification.
\end{block}
\end{frame}

\begin{frame}{Placebos and Negative Controls}
Useful falsification strategies include:
\begin{itemize}
\item placebo treatment dates before the real intervention;
\item outcomes that should not respond to treatment;
\item exposures that should not causally affect the target outcome;
\item re-estimation on untreated comparison units;
\item sensitivity across defensible specifications.
\end{itemize}

\begin{alertblock}{Interpretation}
Passing a placebo check increases credibility only relative to the failure modes that check can detect.
\end{alertblock}
\end{frame}

\section{Applied Workflow}

\begin{frame}{A Reproducible Causal Econometrics Workflow}
\begin{enumerate}
\item Define treatment, outcome, population, timing, and estimand.
\item Write the causal and data-generating story before fitting a model.
\item State identification assumptions explicitly.
\item Choose an estimator implied by the design.
\item Run design-specific diagnostics.
\item Quantify statistical uncertainty.
\item Run falsification and sensitivity analyses.
\item Interpret magnitude, heterogeneity, and external validity.
\end{enumerate}
\end{frame}

\begin{frame}{Example: Regional Pricing Policy}
\textbf{Question:} What is the incremental effect of a regional pricing policy on customer demand?

\begin{columns}[T]
\begin{column}{0.48\textwidth}
\begin{block}{Potential threats}
\begin{itemize}
\item regions selected because demand was already changing;
\item regional macroeconomic shocks;
\item lagged effects;
\item cross-region spillovers;
\item prices respond to expected demand.
\end{itemize}
\end{block}
\end{column}
\begin{column}{0.48\textwidth}
\begin{block}{Candidate designs}
\begin{itemize}
\item fixed effects with rich time controls;
\item DiD and event studies;
\item synthetic control for few treated regions;
\item IV with a defensible exogenous price shifter.
\end{itemize}
\end{block}
\end{column}
\end{columns}
\end{frame}

\begin{frame}{Synthesis: Identification Before Estimation}
\begin{itemize}
\item Prediction asks what will happen; causal econometrics asks what would change under intervention.
\item Fix the estimand before the estimator.
\item Endogeneity is a failure of exogeneity, not merely a modeling inconvenience.
\item Panel FE, DiD, IV, weighting, and synthetic control exploit different identifying variation.
\item Every method inherits assumptions that need substantive defense and diagnostics.
\item Confidence intervals do not protect against identification failure.
\item Strong causal work makes assumptions, failure modes, and sensitivity visible.
\end{itemize}
\end{frame}

\begin{frame}{Selected References}
\small
\begin{itemize}
\item Angrist, J. D., and Pischke, J.-S. (2009). \textit{Mostly Harmless Econometrics}. Princeton University Press.
\item Cunningham, S. (2021). \textit{Causal Inference: The Mixtape}. Yale University Press.
\item Hernan, M. A., and Robins, J. M. (2020). \textit{Causal Inference: What If}. Chapman \& Hall/CRC.
\item Imbens, G. W., and Rubin, D. B. (2015). \textit{Causal Inference for Statistics, Social, and Biomedical Sciences}. Cambridge University Press.
\item Wooldridge, J. M. (2010). \textit{Econometric Analysis of Cross Section and Panel Data}, 2nd ed. MIT Press.
\item Abadie, A. (2021). Using Synthetic Controls: Feasibility, Data Requirements, and Methodological Aspects. \textit{Journal of Economic Literature}, 59(2), 391--425.
\end{itemize}
\end{frame}

\contactslide

\end{document}
